\section{One policy on an arbitrary metric space}
\label{sec:transfer}

Proposition~\ref{prop:finite-law} produces a policy for each input law on
a fixed finite metric.  We now remove that dependence in three steps.
First, repeated independent episodes remove the additive term when the
starting points are distinct.  A finite separation argument then gives
one policy for any finite collection of inputs.  Finally, product compactness gives one policy for all finite inputs,
even when the metric space is infinite. We spell out the finite
separation argument and the closed constraints used in compactness.

Throughout this section, let $a$ be the absolute constant in
Proposition~\ref{prop:finite-law}, and put
\begin{equation}
 \ell=1+\log(k+1),\qquad c_{\mathrm{tr}}=2a+1,
 \qquad \alpha=c_{\mathrm{tr}}\ell^2.
 \label{eq:transfer-alpha}
\end{equation}
We may enlarge the additive constant of that proposition to ensure
$D_0(Y,k,s)\ge0$.  Its independence of the input law and horizon is
essential below.

\subsection{Removing the finite-metric additive constant}

\begin{lemma}[Episode amplification]
\label{lem:amplification}
Let $Y$ be a finite metric space, let $k\ge2$, and let
$s=(s_1,\ldots,s_k)\in Y^k$ have distinct entries.  For every
finite-support probability law $\lambda$ on finite request sequences in
$Y$, there is a causal randomized policy $A_\lambda$ starting from $s$
such that
\begin{equation}
 \E_{\sigma\sim\lambda}\E\bigl[\cost_{A_\lambda,s}(\sigma)\bigr]
 \le \alpha\,
      \E_{\sigma\sim\lambda}\bigl[\OPT_s(\sigma)\bigr].
 \label{eq:additive-free-law}
\end{equation}
The policy can be strongly lazy.  The constant $c_{\mathrm{tr}}$ in
\eqref{eq:transfer-alpha} is the same for all $Y$, $s$, and $\lambda$.
\end{lemma}

\begin{proof}
Write $Q=\E_{\lambda}\OPT_s(\sigma)$.  We first make the episode length
deterministic without changing $Q$.  Let
\[
 L=\max\bigl(1,\max_{\sigma\in\supp\lambda}|\sigma|\bigr).
\]
Extend each nonempty sequence to length $L$ by repeating its last
request; extend an empty sequence by $L$ repetitions of $s_1$.
Denote the resulting law on $Y^L$ by $\bar\lambda$.  For every original
sequence $\sigma$ and its extension $\bar\sigma$,
\begin{equation}
 \OPT_s(\bar\sigma)=\OPT_s(\sigma).
 \label{eq:padding-optimum}
\end{equation}
Indeed, restriction of any service of $\bar\sigma$ to its original
prefix gives the lower bound.  An optimal service of a nonempty
$\sigma$ ends with a server at its last request and can serve the
extension at zero extra cost.  An empty sequence and its extension both
have optimum zero.  A causal policy designed for $L$ requests can be
run on the original sequence and stopped when that sequence ends; its
actions on this prefix are identical to its actions on the extension.
Thus it is enough to prove the claim for $\bar\lambda$ and then
truncate.  The policy need not be told the actual stopping time in
advance.

If $Q=0$, every sequence of positive probability has optimum zero.
An optimum is attained because there are only $k^L$ possible strings of
server labels.  A zero-cost service cannot move any server to a
different point, so every request lies among $s_1,\ldots,s_k$.
Any strongly lazy policy therefore has zero cost on the law, proving
\eqref{eq:additive-free-law}.  We henceforth assume $Q>0$.

Set
\[
 \Delta=\diam(Y),\qquad
 \delta=\min\{d(y,z):y,z\in Y,\ y\ne z\}>0,
 \qquad R=\left\lfloor\frac{k\Delta}{\delta}\right\rfloor+1.
\]
The distinct starting points ensure $\Delta>0$.  The reset word is
$R$ consecutive copies of $(s_1,\ldots,s_k)$, of length $Rk$.
For every $\rho=(x_1,\ldots,x_L)\in\supp\bar\lambda$, choose one
optimal labeled service from $s$, resolving ties by the lexicographic
order of its chosen server labels.  Write $C_t(\rho)\in Y^k$ for its
configuration after request $t$, and $j_t(\rho)$ for its chosen label.
Here $C_0(\rho)=s$; a label is designated even when the move is zero.
Define its reverse word to be
\begin{equation}
 \rho^{\leftarrow}
  =\bigl(C_{L-1,j_L}(\rho),C_{L-2,j_{L-1}}(\rho),
              \ldots,C_{0,j_1}(\rho)\bigr).
 \label{eq:reverse-word}
\end{equation}
Starting at $C_L(\rho)$, this word can be served by moving labels
$j_L,j_{L-1},\ldots,j_1$ in that order.  Before undoing step $t$ the
configuration is exactly $C_t(\rho)$, and its designated move returns
it to $C_{t-1}(\rho)$.  This verifies legality and returns the labeled
tuple exactly to $s$, at total cost $\OPT_s(\rho)$.  The assertion also
covers zero moves.  The reset word then has zero comparator cost.

Fix an integer $N\ge1$.  Sample $N$ independent words
$\rho_1,\ldots,\rho_N$ from $\bar\lambda$, and concatenate, for each
$i$, its forward word $\rho_i$, its reverse word
$\rho_i^{\leftarrow}$, and the reset word.  This is a finite-support
law on words of the fixed length $N(2L+Rk)$.  All its requests are
functions of the sampled forward words and the fixed optimal services;
in particular, the law is fixed independently of the online policy's
randomness.  Following the chosen optimal service and then its reverse
in every episode gives
\begin{equation}
 \E\OPT_s(\text{concatenated experiment})\le 2NQ.
 \label{eq:episode-comparator}
\end{equation}
Proposition~\ref{prop:finite-law}, followed if necessary by
Lemma~\ref{lem:lazy-simulation} with zero initial mismatch, gives a
strongly lazy causal policy $\mathcal A_N$ for this experiment with
\begin{equation}
 \E\cost_{\mathcal A_N,s}(\text{experiment})
       \le 2a\ell^2NQ+D_0(Y,k,s).
 \label{eq:episode-total}
\end{equation}
As usual, specify arbitrary causal strongly lazy continuation rules at
histories having probability zero under the experiment law.

Let $F_i$, $H_i$, and $Z_i$ be the costs of $\mathcal A_N$ during the
forward word, reverse word, and reset word of episode $i$,
respectively.  Let $A_i$ be its configuration at the end of that reset
word, and define
\[
 M_i=\mathsf M(A_i,s),\qquad M_0=0.
\]
We claim, pathwise, that
\begin{equation}
 M_i\le Z_i\qquad(1\le i\le N).
 \label{eq:reset-mismatch}
\end{equation}
If a complete reset cycle has zero movement, the configuration stays
constant throughout that cycle and covers every one of the $k$
distinct points $s_1,\ldots,s_k$.  Since there are exactly $k$ servers,
its multiset of positions is the starting multiset.  Strong laziness
keeps this multiset fixed throughout all remaining reset cycles.
Consequently, if $M_i>0$, every reset cycle contains a positive move.
Each such move has length at least $\delta$, because all positions
are in $Y$.  Therefore
\[
 Z_i\ge R\delta>k\Delta\ge\mathsf M(A_i,s)=M_i.
\]
If $M_i=0$, \eqref{eq:reset-mismatch} follows from nonnegativity.
The argument requires equality only of multisets, not of the labels
occupying the distinct starting points.

We extract a policy for one forward word.  Before receiving any actual
requests, choose $I$ uniformly from $\{1,\ldots,N\}$.  Conditional on
$I=i$, generate $i-1$ fictitious independent forward words from
$\bar\lambda$, together with their fixed reverse and reset words.
Using fresh internal randomness independent of these words, simulate
$\mathcal A_N$ from $s$ through that entire fictitious prefix.  This
produces its internal state and a virtual configuration at the start
of forward episode $i$.  For each actual request, now give that
request to the virtual policy as its next request.  Follow its virtual
configurations from the real initial tuple $s$ by the strongly lazy
simulation in Lemma~\ref{lem:lazy-simulation}.

This defines a causal policy: the fictitious prefix is internal
randomness independent of the actual input, the boundary of episode
$i$ is known from $i$, $L$, and $R$, and the virtual continuation uses
only the actual requests already received.  Neither the remainder of
the actual word nor its reverse word needs to be generated.  After
$L$ actual requests, an arbitrary strongly lazy rule extends this
policy to all longer inputs.

When the actual forward word is sampled from $\bar\lambda$,
conditional on $I=i$, the virtual evolution through that word has the
same law as the evolution of $\mathcal A_N$ through forward episode
$i$ in the experiment.  The true simulation starts with mismatch
$M_{i-1}$ under this identification, where $M_0=0$ for $i=1$.
Lemma~\ref{lem:lazy-simulation} and \eqref{eq:reset-mismatch} give
\begin{align}
 \E\cost_{\mathrm{extracted},s}(\rho)
 &\le \frac1N\sum_{i=1}^N\E(F_i+M_{i-1}) \notag\\
 &\le \frac1N\left(\sum_{i=1}^N\E F_i
                    +\sum_{i=1}^{N-1}\E Z_i\right) \notag\\
 &\le \frac1N\E\cost_{\mathcal A_N,s}(\text{experiment}) \notag\\
 &\le 2a\ell^2Q+\frac{D_0(Y,k,s)}N.
 \label{eq:finite-n-amplification}
\end{align}
In the third line, the charged forward and preceding reset blocks are
disjoint parts of the experiment; reverse-block costs and the last
reset-block cost are nonnegative uncharged terms.  Truncating to an
original unpadded sequence can only reduce this expected cost.

Finally take
\[
 N=\max\left(1,
      \left\lceil\frac{D_0(Y,k,s)}{\ell^2Q}\right\rceil\right).
\]
Then \eqref{eq:finite-n-amplification} is at most
$(2a+1)\ell^2Q$, as asserted.  This is a finite choice of $N$ for the
given law; no limiting policy is used here.  The metric-dependent
numbers $R$ and $N$ affect only the auxiliary experiment and not the
coefficient in \eqref{eq:additive-free-law}.
\end{proof}

\subsection{A single policy for a finite collection of inputs}

For this and the next subsection fix an arbitrary metric space $X$
and a distinct starting tuple $s\in X^k$.  Write $[k]=\{1,\ldots,k\}$.
A decision history at time $t$ consists of the request prefix
$(x_1,\ldots,x_t)$ and the previously chosen server labels
$(j_1,\ldots,j_{t-1})$.  These data determine the current labeled
configuration.  A behavioral policy assigns a probability vector on
$[k]$ to every such history and samples its next label from that vector.
It has no horizon parameter.

\begin{lemma}[Finite separation and behavioral conversion]
\label{lem:finite-minimax}
For every finite collection $\sigma_1,\ldots,\sigma_m$ of finite
request sequences in $X$, there is a behavioral policy $p$ from $s$
satisfying
\begin{equation}
 \E\cost_{p,s}(\sigma_i)\le\alpha\OPT_s(\sigma_i)
       \qquad(1\le i\le m).
 \label{eq:finite-family-bound}
\end{equation}
Its behavioral choices can be defined on all decision histories in $X$.
\end{lemma}

\begin{proof}
There is nothing to prove for an empty collection.  Otherwise let $Y$
be the finite submetric containing all starting points and all points
appearing in the sequences.  It has positive diameter.  Every server
position during a legal service is an initial position or a previously
requested point: each permitted move is to the current request.
Therefore passing from $X$ to $Y$ changes neither the legal label
strings nor their costs, and
\begin{equation}
 \OPT_{Y,s}(\sigma_i)=\OPT_{X,s}(\sigma_i).
 \label{eq:finite-restriction-optimum}
\end{equation}

Let $\mathcal H_0$ be the finite set of decision histories whose request
prefix is a nonempty prefix of some $\sigma_i$, with every possible
previous label string included.  This is the finite distributional minimax principle associated with
Yao \cite{Yao77}, applied to excess cost. A deterministic policy on this set is
a map $\pi:\mathcal H_0\to[k]$.  There are finitely many such maps;
write their set as $\mathcal D$.  For each $\pi\in\mathcal D$, put
\[
 v(\pi)_i=\cost_{\pi,s}(\sigma_i)-\alpha\OPT_s(\sigma_i),
 \qquad
 K=\left\{\sum_{\pi\in\mathcal D}\theta_\pi v(\pi):
       \theta_\pi\ge0,\ \sum_{\pi\in\mathcal D}\theta_\pi=1\right\}.
\]
The set $K\subseteq\R^m$ is compact, being the continuous image of a
finite-dimensional compact simplex, and is convex.

We claim $K\cap(-\R_{\ge0}^m)\ne\varnothing$.  If not, minimize the
continuous function
\[
 f(z)=\sum_{i=1}^m\max(z_i,0)^2
\]
over $K$, and let $z_*$ be a minimizer.  Put
$w_i=\max((z_*)_i,0)$.  The supposed disjointness implies $w\ne0$.
For every $z\in K$ the segment $z_*+t(z-z_*)$, $0\le t\le1$, lies in
$K$.  The right derivative of $f$ along this segment at $t=0$ is
nonnegative.  Since the scalar function $u\mapsto\max(u,0)^2$ is
differentiable, including at zero, this says
\[
 2w\mathbin{\cdot}(z-z_*)\ge0.
\]
It follows that
\begin{equation}
 w\mathbin{\cdot}z\ge w\mathbin{\cdot}z_*
    =\sum_i w_i^2>0\qquad(z\in K).
 \label{eq:finite-separation}
\end{equation}
Normalize $w$ to the probability weights
$\lambda_i=w_i/(\sum_jw_j)$ on the target sequences.
Lemma~\ref{lem:amplification}, applied on $Y$, supplies a randomized
causal policy with nonpositive expected excess cost for this law.

Every randomized causal policy, restricted to this finite history set,
has its vector of expected excess costs in $K$.  To see this directly,
fix its internal random tape, independent of requests.  Run the fixed
tape on each of the finitely many request prefixes.  This determines
a deterministic rule at histories compatible with those runs; assign
arbitrary choices at incompatible histories.  The resulting member of
$\mathcal D$ reproduces its actions on every $\sigma_i$.
There are only finitely many possible rules, so their probabilities
form a convex combination of the vectors $v(\pi)$.  Each event defining
such a rule is a finite combination of events for choices on fixed
finite inputs and is measurable.  Consequently the expected excess
for the law with weights $\lambda_i$ is strictly positive by
\eqref{eq:finite-separation}, a contradiction.  This proves the claim.
Choose weights $\theta_\pi$ whose convex combination lies in the
nonpositive orthant.  Sampling $\pi$ with these weights at the start
would give \eqref{eq:finite-family-bound}.

For later compactness we convert this finite mixture to behavioral
choices, without placing any law on requests. This is the
perfect-recall conversion of mixed to behavioral strategies
\cite{Kuhn53}; the formula below makes its request-independent
meaning explicit.  At a history
$h=((x_1,\ldots,x_t),(j_1,\ldots,j_{t-1}))\in\mathcal H_0$, let
\[
 \mathcal D(h)=\{\pi\in\mathcal D:
   \pi((x_1,\ldots,x_u),(j_1,\ldots,j_{u-1}))=j_u
                    \text{ for every }u<t\},
 \qquad Z(h)=\sum_{\pi\in\mathcal D(h)}\theta_\pi.
\]
If $Z(h)>0$, define
\begin{equation}
 p_h(j)=\frac{\displaystyle
          \sum_{\substack{\pi\in\mathcal D(h)\\\pi(h)=j}}
               \theta_\pi}{Z(h)}\qquad(j\in[k]).
 \label{eq:behavioral-conversion}
\end{equation}
If $Z(h)=0$, choose any probability vector, for example one concentrated
on label $1$.  Induction along a fixed request sequence shows that the
probability of each label prefix under these choices is the total
$\theta$-weight of policies consistent with that prefix.  Indeed, when
this weight is positive the next factor in
\eqref{eq:behavioral-conversion} is exactly the ratio of the extended
prefix weight to the old weight; when it is zero the prefix is
unreachable.  Thus the complete label strings, and therefore their
costs, have the same law as in the original mixture on each $\sigma_i$.
Set $p_h$ arbitrarily at histories outside $\mathcal H_0$ to obtain the
asserted policy on all of $X$.  Formula~\eqref{eq:behavioral-conversion}
conditions only on consistency of past choices along a specified
request prefix, not on an assumed distribution of requests.
\end{proof}

\subsection{Compactness over all finite histories}

\begin{lemma}[One policy for all finite inputs]
\label{lem:compact-policy}
For every metric space $X$ and every distinct tuple $s\in X^k$, there
is a single randomized causal policy $A_s$ such that
\begin{equation}
 \E\cost_{A_s,s}(\sigma)\le\alpha\OPT_s(\sigma)
       \qquad\text{for every finite request sequence }\sigma\text{ in }X.
 \label{eq:all-inputs-distinct}
\end{equation}
For each fixed input its decisions and cost are measurable functions
of its internal randomness.  Neither its rule nor its randomization
depends on an eventual horizon.
\end{lemma}

\begin{proof}
The set of all decision histories is
\[
 \mathcal H=\bigcup_{t\ge1}\bigl(X^t\times[k]^{t-1}\bigr).
\]
Let
\[
 \Delta_k=\{(p(1),\ldots,p(k))\in[0,1]^k:
                       \textstyle\sum_{j=1}^kp(j)=1\},
 \qquad \mathcal P=\prod_{h\in\mathcal H}\Delta_k.
\]
Thus a point of $\mathcal P$ specifies a behavioral policy at all
histories. Each factor $\Delta_k$ is compact, so Tychonoff's theorem
makes $\mathcal P$ compact in the product topology
\cite[Chapter~5, Theorem~13]{Kelley55}; see also \cite{Tychonoff35}.
This application uses the usual set-theoretic
choice principle and permits an arbitrary index set $\mathcal H$.

First record the exact dependence of a fixed-input expected cost on
these coordinates.  For $\sigma=(x_1,\ldots,x_m)$ and a label string
$j=(j_1,\ldots,j_m)\in[k]^m$, let $c_s(\sigma,j)$ be its deterministic
service cost from $s$.  Explicitly, begin with the tuple $s$, replace
coordinate $j_t$ by $x_t$ at step $t$, and sum the distances replaced.
These costs are finite, even when $X$ is unbounded, since a fixed input
and the starting tuple contain only finitely many points.  For
$p\in\mathcal P$ the expected cost is
\begin{equation}
 C_\sigma(p)=\sum_{j\in[k]^m}c_s(\sigma,j)
       \prod_{t=1}^m
       p_{((x_1,\ldots,x_t),(j_1,\ldots,j_{t-1}))}(j_t).
 \label{eq:cost-polynomial}
\end{equation}
For the empty sequence this is zero.  For any fixed nonempty sequence,
it is a polynomial in finitely many coordinates of $p$, hence a
continuous function of those coordinates.  Consequently the constraint
\begin{equation}
 C_\sigma(p)\le\alpha\OPT_s(\sigma)
 \label{eq:closed-cost-constraint}
\end{equation}
is closed in those finitely many factors.

For each finite input $\sigma$, let
\[
 E_\sigma=\{p\in\mathcal P:
 C_\sigma(p)\le\alpha\OPT_s(\sigma)\}.
\]
These sets are closed by the finite-coordinate continuity just proved.
Every finite subcollection has nonempty intersection by
Lemma~\ref{lem:finite-minimax}: that lemma supplies a behavioral
policy defined on all histories which satisfies all the selected
constraints. The finite-intersection property in the compact space
$\mathcal P$ therefore gives
\[
                         \bigcap_\sigma E_\sigma\ne\varnothing.
\]
Choose $p$ in this intersection. No bound on the diameter or cardinality
of $X$ enters this argument. Boundedness is needed only for the cost of
one fixed finite input, which was established before
Equation~\eqref{eq:cost-polynomial}.

It remains to realize these choices as one randomized online policy.
At each request use the current history $h$ and choose the next label
with probabilities $p_h(1),\ldots,p_h(k)$.  For example, take
independent uniform random numbers in $[0,1)$ and choose the first
label whose cumulative probability exceeds the next number.  Only the
revealed request prefix and past labels enter this rule.
One can also implement all times on the single probability space
$[0,1)$ with Lebesgue measure: subdivide its current interval into
half-open pieces with relative lengths $p_h(1),\ldots,p_h(k)$, choose
the piece containing the initial uniform point, and continue within
that piece.  Pieces of length zero are empty.  On a fixed finite input,
the set of seeds giving any specified label prefix is an interval
whose length is the product of the corresponding probabilities.
Therefore each label choice and the finite total cost are measurable,
and their expectation is precisely \eqref{eq:cost-polynomial}.

In particular, the policy is defined for every history before any
input or horizon is chosen.  Measurability here is exactly that required
for each fixed oblivious input; no probability distribution on the
request space, or joint measurability in a random request variable,
enters the construction.  The inequalities already proved for $p$
give \eqref{eq:all-inputs-distinct}.
\end{proof}

\subsection{Repeated starting positions and the final bound}

\begin{proof}[Proof of Theorem~\ref{thm:main}]
Fix $k\ge2$, a metric space $(X,d)$ with at least $k+1$ points, and any
labeled starting tuple $s\in X^k$, allowing repetitions.  Choose a
tuple $u\in X^k$ of distinct points, and put
\begin{equation}
 D=\mathsf M(s,u)
   =\min_{\pi\text{ a permutation of }[k]}
          \sum_{i=1}^k d(s_i,u_{\pi(i)}).
 \label{eq:initial-matching-distance}
\end{equation}
Such a tuple exists by the cardinality assumption.  The number $D$ is
finite, even in an unbounded metric, since it is the minimum of
finitely many finite sums.

For every finite input $\sigma$, choose an optimal legal trajectory
from $s$.  Its existence follows from the finitely many label strings
for that input.  Apply Lemma~\ref{lem:lazy-simulation} to follow this
trajectory from the initial tuple $u$.  This gives a legal service from
$u$ of cost at most $\OPT_s(\sigma)+D$.  This use of the lemma is a
pathwise comparison of offline services and does not require computing
that trajectory online.  Consequently
\begin{equation}
 \OPT_u(\sigma)\le\OPT_s(\sigma)+D.
 \label{eq:offline-start-comparison}
\end{equation}

By Lemma~\ref{lem:compact-policy}, one online policy $A_u$ from $u$
satisfies \eqref{eq:all-inputs-distinct} for every finite input.
Run this policy virtually, and follow its configurations from $s$ by
Lemma~\ref{lem:lazy-simulation}.  The resulting policy $A$ is causal
and can be strongly lazy.  Minimum matchings can be selected by a fixed
order on the finite set of label permutations, so this simulation also
preserves measurability on each fixed input.  Pathwise it adds at most
$D$ to the virtual cost, and therefore
\begin{align}
 \E\cost_{A,s}(\sigma)
 &\le \E\cost_{A_u,u}(\sigma)+D \notag\\
 &\le \alpha\OPT_u(\sigma)+D \notag\\
 &\le \alpha\OPT_s(\sigma)+(\alpha+1)D.
 \label{eq:repeated-start-bound}
\end{align}
The explicit additive constant
\begin{equation}
 B=\bigl((2a+1)(1+\log(k+1))^2+1\bigr)\mathsf M(s,u)
 \label{eq:final-additive-constant}
\end{equation}
depends only on the metric, $k$, and the fixed choices of starting
tuples.  It is independent of the input and its length.  When $s$
already has distinct entries we may take $u=s$, giving $D=B=0$.

Finally, for $k\ge2$ we have $\log(k+1)\ge\log3>1$ and hence
\[
 (1+\log(k+1))^2\le4(\log(k+1))^2.
\]
Thus \eqref{eq:repeated-start-bound} proves the theorem with the single
absolute constant $C=4(2a+1)$ and the finite $B$ in
\eqref{eq:final-additive-constant}.  All policy choices were fixed for
the specified metric and initial tuple, so the same policy serves
every finite sequence and every horizon.
\end{proof}
